% ============================================================
% Results section — EMNLP paper
% ============================================================

\section{Results}

\subsection{Main Results}

Table~\ref{tab:main_results} reports Weighted F1 (W-F1), location accuracy (Loc), and
joint accuracy (Joint) for all models under the Baseline and our full method (+GI+SI)
on TRAIL-GAIA and TRAIL-SWE-Bench.
Table~\ref{tab:trace_coverage} reports the number of traces actually processed ($N$)
per model, dataset, and method; we note that $N$ is not always identical across
conditions for the same model, which we discuss below.

\paragraph{TRAIL-GAIA.}
Augmenting agent traces with causal graph injection and span-index prefixes (+GI+SI)
consistently improves over the no-graph baseline across all six models on TRAIL-GAIA.
W-F1 gains range from $+3.7$ points (Gemini-2.5-Flash) to $+9.1$ points
(Gemma-3-27B-IT), demonstrating that causal context extracted from execution traces
provides a reliable signal for error categorization regardless of model family or
scale.
Crucially, even the strongest evaluated model, Gemini-2.5-Flash, benefits from the
augmentation, ruling out the interpretation that gains arise solely from weaker models
relying on graph scaffolding as a shortcut.

Location accuracy improves yet more dramatically in several cases.
QwenLong-L1-32B shows a $+15.3$-point Loc gain ($3.92 \to 19.23$) and
GPT-oss-20B improves by $+12.0$ points ($6.29 \to 18.29$).
Notably, GPT-oss-20B simultaneously exhibits a marginal W-F1 regression ($-1.0$),
suggesting a dissociation between the two sub-tasks: span-index prefixes effectively
guide the model toward the correct fault-bearing step, but correctly labeling the
error category at that step remains challenging for smaller models.
When both sub-tasks succeed together, the gains compound: Mistral-Small-3.1-24B
achieves the largest absolute joint-accuracy improvement ($+7.7$ points,
$3.78 \to 11.51$), and QwenLong-L1-32B gains $+3.4$ joint points.

\paragraph{Context window as a first-order constraint.}
TRAIL traces can be extremely long — full-trace inputs routinely exceed
100K tokens — and the additional token overhead introduced by the span index and
causal graph description interacts critically with each model's available context
budget.
This constraint shapes which conditions can be reliably evaluated.
For QwenLong-L1-32B on the original (non-dedup) full-length traces, adding the
static causal graph together with the span index (+CG+SI) triggers complete
context overflow, producing \textbf{zero valid predictions} ($\text{W-F1} = 0.00$).
The evaluator fix that counts context-length truncations as empty predictions
rather than skipping them~(§\ref{sec:eval}) makes such failures visible rather
than silently inflating scores.
The trace-coverage data in Table~\ref{tab:trace_coverage} reveal a complementary
signal: QwenLong's SWE +CG+SI condition has the lowest processed-trace count
across all experiments ($N{=}16$), and Mistral's +GI+SI run on GAIA covers only
64 of the available 81 traces, both pointing to runs that stalled under token-budget
pressure before completing the full evaluation set.

The \emph{design} of +GI+SI directly addresses this constraint.
Unlike +CG+SI, which prepends the full edge list and span index to a single
monolithic prompt, +GI+SI runs a short, targeted second pass containing only the
1–3 downstream error categories predicted from the first pass.
This limits the token overhead of the causal hints to the subset of edges actually
activated by what the model detected, keeping the augmented prompt within budget
even for models where a static full-graph prompt would overflow.
The token efficiency of the two-pass design therefore contributes directly to
+GI+SI's superior aggregate performance, particularly for models operating near
their effective context limits.

\paragraph{TRAIL-SWE-Bench.}
SWE-Bench traces present a substantially harder fault-localization challenge than GAIA.
Baseline W-F1 scores are far lower for all models (e.g., Gemini-2.5-Flash: $8.71$
vs.\ $37.07$ on GAIA), and location accuracy remains near zero across the board even
after augmentation.
A contributing factor to the difficulty is trace coverage: SWE $N$ ranges from only
$9$ to $22$ across all models and conditions (Table~\ref{tab:trace_coverage}),
roughly one-fifth the scale of the GAIA runs, so individual SWE scores carry wider
confidence intervals and should be treated as indicative rather than definitive.
This systematic gap is at least partly attributable to trace length: SWE-Bench traces
capture repository-level software engineering sessions with extended code execution
histories, which are longer and denser than the multi-step question-answering traces
in GAIA.
Under tighter effective context budgets, the fault-bearing steps — which often occur
early in the agent's session, where the root cause is introduced — are disproportionately
likely to be truncated before the model processes the full trace, making precise
localization structurally harder regardless of the augmentation method.

Despite this, +GI+SI yields positive W-F1 gains for five of six models on SWE-Bench.
Gemini-2.5-Flash shows the largest absolute improvement ($+12.9$ W-F1,
$8.71 \to 21.56$), followed by QwenLong-L1-32B ($+9.0$ W-F1, $2.14 \to 11.09$).
Mistral-Small-3.1-24B is the sole exception ($-3.0$ W-F1), suggesting that for models
with limited code-reasoning capacity and no compressed-trace fallback on the SWE task,
the injected graph structure may consume context budget that would otherwise be used
for trace content, introducing noise rather than signal.
Joint accuracy is at or near zero for all models on SWE-Bench, underscoring that
simultaneous step localization and error categorization in code-oriented traces
remains an open challenge.

\subsection{Ablation Studies}

\paragraph{Graph augmentation variants.}
Table~\ref{tab:ablation_methods} ablates the four method variants on TRAIL-GAIA using
Gemini-2.5-Flash, whose large context window eliminates context overflow as a
confound, isolating the contribution of each design choice.
All graph-augmented conditions outperform the baseline, confirming that causal
structure is beneficial regardless of presentation.
Adding only the causal graph (+CG) already recovers most of the W-F1 gain
($40.21$ vs.\ $40.75$ for +GI+SI), indicating that graph topology itself is the
primary source of diagnostic signal.
Adding the span index alongside the static graph (+CG+SI) yields the highest
joint accuracy ($15.81$, $+3.1$ over baseline), as the closed-vocabulary span index
converts the location field from a free-form generation problem into a
selection problem, suppressing span ID hallucination.
The full +GI+SI method achieves the best W-F1 ($40.75$) while matching +CG on
location accuracy ($35.14$), with the conditional two-pass design suppressing the
over-triggering that causes +CG to propose spurious downstream error categories
when the evidence is weak.
For models without Gemini's context headroom, this conditionality also provides
the token-budget advantage described above, making +GI+SI the most balanced
configuration across scales.

\paragraph{Effect of dataset deduplication.}
Table~\ref{tab:ablation_dedup} compares Mistral-Small-3.1-24B across the original and
deduplicated TRAIL-GAIA splits.
On the original split, the model achieves only $17.98$ W-F1 under +CG and $20.85$
under +CG+SI; no clean baseline is available due to run constraints.
On the deduplicated split, the same graph conditions yield substantially higher scores:
+CG reaches $28.31$ W-F1 ($+10.3$ absolute over the original split), and +GI+SI
reaches $30.76$, with a clean baseline of $24.06$ confirming a true per-condition
comparison.
The original split contains near-duplicate instances that both inflate apparent baselines
and compress per-condition variance, obscuring the true contribution of graph augmentation.
Notably, context-length truncations — which were silently skipped in earlier evaluator
versions — are concentrated in the longer traces that near-duplicate instances tend to
produce, creating an additional systematic bias in original-split scores.
Consequently, all small-model evaluations are reported on the deduplicated splits to
ensure reliable, contamination-free measurement.
